\documentclass[10pt,a4paper]{amsart}
\usepackage{amsmath,amssymb,amsthm,graphicx}
% For commutative diagrams we use Xy-pic
\usepackage[all]{xy}

% We use 2cell for 2-commutative diagrams.
\xyoption{2cell}
% Theorem environments.
%
\theoremstyle{plain}
\newtheorem{theorem}[subsection]{Theorem}
\newtheorem{proposition}[subsection]{Proposition}
\newtheorem{lemma}[subsection]{Lemma}

\theoremstyle{definition}
\newtheorem{definition}[subsection]{Definition}
\newtheorem{example}[subsection]{Example}
\newtheorem{exercise}[subsection]{Exercise}
\newtheorem{situation}[subsection]{Situation}

\theoremstyle{remark}
\newtheorem{remark}[subsection]{Remark}
\newtheorem{remarks}[subsection]{Remarks}

\numberwithin{equation}{subsection}

% Macros
%
\def\lim{\mathop{\mathrm{lim}}\nolimits}
\def\colim{\mathop{\mathrm{colim}}\nolimits}
\def\Spec{\mathop{\mathrm{Spec}}}
\def\Hom{\mathop{\mathrm{Hom}}\nolimits}
\def\Ext{\mathop{\mathrm{Ext}}\nolimits}
\def\SheafHom{\mathop{\mathcal{H}\!\mathit{om}}\nolimits}
\def\SheafExt{\mathop{\mathcal{E}\!\mathit{xt}}\nolimits}
\def\Sch{\mathit{Sch}}
\def\Mor{\mathop{\mathrm{Mor}}\nolimits}
\def\Ob{\mathop{\mathrm{Ob}}\nolimits}
\def\Sh{\mathop{\mathit{Sh}}\nolimits}
\def\NL{\mathop{N\!L}\nolimits}
\def\CH{\mathop{\mathrm{CH}}\nolimits}
\def\proetale{{pro\text{-}\acute{e}tale}}
\def\etale{{\acute{e}tale}}
\def\QCoh{\mathit{QCoh}}
\def\Ker{\mathop{\mathrm{Ker}}}
\def\Im{\mathop{\mathrm{Im}}}
\def\Coker{\mathop{\mathrm{Coker}}}
\def\Coim{\mathop{\mathrm{Coim}}}

% Boxtimes
%
\DeclareMathSymbol{\boxtimes}{\mathbin}{AMSa}{"02}

%
% Macros for moduli stacks/spaces
%
\def\QCohstack{\mathcal{QC}\!\mathit{oh}}
\def\Cohstack{\mathcal{C}\!\mathit{oh}}
\def\Spacesstack{\mathcal{S}\!\mathit{paces}}
\def\Quotfunctor{\mathrm{Quot}}
\def\Hilbfunctor{\mathrm{Hilb}}
\def\Curvesstack{\mathcal{C}\!\mathit{urves}}
\def\Polarizedstack{\mathcal{P}\!\mathit{olarized}}
\def\Complexesstack{\mathcal{C}\!\mathit{omplexes}}
% \Pic is the operator that assigns to X its picard group, usage \Pic(X)
% \Picardstack_{X/B} denotes the Picard stack of X over B
% \Picardfunctor_{X/B} denotes the Picard functor of X over B
\def\Pic{\mathop{\mathrm{Pic}}\nolimits}
\def\Picardstack{\mathcal{P}\!\mathit{ic}}
\def\Picardfunctor{\mathrm{Pic}}
\def\Deformationcategory{\mathcal{D}\!\mathit{ef}}
\usepackage{cleveref}
\crefname{originallemma}{original lemma}{original lemmas}
\Crefname{originallemma}{Original lemma}{Original lemmas}
\makeatletter
\newlabel{homology-lemma-characterize-weak-serre-subcategory}{{0754}{source}{}{sourcepointer-1}{}}
\newlabel{homology-lemma-characterize-weak-serre-subcategory@cref}{{[originallemma][0][]0754}{[1][0][]source}{}{}{}}
\makeatother
\crefname{originalreference}{original reference}{original references}
\Crefname{originalreference}{Original reference}{Original references}
\makeatletter
\newlabel{lemma-cohomological-descent-etale}{{0DDI}{source}{}{sourcepointer-2}{}}
\newlabel{lemma-cohomological-descent-etale@cref}{{[originalreference][0][]0DDI}{[1][0][]source}{}{}{}}
\makeatother
\crefname{originalreference}{original reference}{original references}
\Crefname{originalreference}{Original reference}{Original references}
\makeatletter
\newlabel{lemma-compare-higher-direct-image-proper}{{0DDJ}{source}{}{sourcepointer-3}{}}
\newlabel{lemma-compare-higher-direct-image-proper@cref}{{[originalreference][0][]0DDJ}{[1][0][]source}{}{}{}}
\makeatother
\crefname{originalreference}{original reference}{original references}
\Crefname{originalreference}{Original reference}{Original references}
\makeatletter
\newlabel{lemma-describe-pullback-pi-ph}{{0DDW}{source}{}{sourcepointer-4}{}}
\newlabel{lemma-describe-pullback-pi-ph@cref}{{[originalreference][0][]0DDW}{[1][0][]source}{}{}{}}
\makeatother
\crefname{originalreference}{original reference}{original references}
\Crefname{originalreference}{Original reference}{Original references}
\makeatletter
\newlabel{lemma-torsion-direct-image}{{0DDD}{source}{}{sourcepointer-5}{}}
\newlabel{lemma-torsion-direct-image@cref}{{[originalreference][0][]0DDD}{[1][0][]source}{}{}{}}
\makeatother
\crefname{originallemma}{original lemma}{original lemmas}
\Crefname{originallemma}{Original lemma}{Original lemmas}
\makeatletter
\newlabel{more-morphisms-lemma-dominate-fppf-etale-locally}{{0DET}{source}{}{sourcepointer-6}{}}
\newlabel{more-morphisms-lemma-dominate-fppf-etale-locally@cref}{{[originallemma][0][]0DET}{[1][0][]source}{}{}{}}
\makeatother
\crefname{originallemma}{original lemma}{original lemmas}
\Crefname{originallemma}{Original lemma}{Original lemmas}
\makeatletter
\newlabel{morphisms-lemma-base-change-proper}{{01W4}{source}{}{sourcepointer-7}{}}
\newlabel{morphisms-lemma-base-change-proper@cref}{{[originallemma][0][]01W4}{[1][0][]source}{}{}{}}
\makeatother
\crefname{originalreference}{original reference}{original references}
\Crefname{originalreference}{Original reference}{Original references}
\makeatletter
\newlabel{sites-cohomology-item-A-and-P}{{0EZ7}{source}{}{sourcepointer-8}{}}
\newlabel{sites-cohomology-item-A-and-P@cref}{{[originalreference][0][]0EZ7}{[1][0][]source}{}{}{}}
\makeatother
\crefname{originalreference}{original reference}{original references}
\Crefname{originalreference}{Original reference}{Original references}
\makeatletter
\newlabel{sites-cohomology-item-A-sheaf}{{0EZ6}{source}{}{sourcepointer-9}{}}
\newlabel{sites-cohomology-item-A-sheaf@cref}{{[originalreference][0][]0EZ6}{[1][0][]source}{}{}{}}
\makeatother
\crefname{originalreference}{original reference}{original references}
\Crefname{originalreference}{Original reference}{Original references}
\makeatletter
\newlabel{sites-cohomology-item-base-change-P}{{0EZ4}{source}{}{sourcepointer-10}{}}
\newlabel{sites-cohomology-item-base-change-P@cref}{{[originalreference][0][]0EZ4}{[1][0][]source}{}{}{}}
\makeatother
\crefname{originalreference}{original reference}{original references}
\Crefname{originalreference}{Original reference}{Original references}
\makeatletter
\newlabel{sites-cohomology-item-refine-tau-by-P}{{0EZ8}{source}{}{sourcepointer-11}{}}
\newlabel{sites-cohomology-item-refine-tau-by-P@cref}{{[originalreference][0][]0EZ8}{[1][0][]source}{}{}{}}
\makeatother
\crefname{originalreference}{original reference}{original references}
\Crefname{originalreference}{Original reference}{Original references}
\makeatletter
\newlabel{sites-cohomology-item-restriction-A}{{0EZ5}{source}{}{sourcepointer-12}{}}
\newlabel{sites-cohomology-item-restriction-A@cref}{{[originalreference][0][]0EZ5}{[1][0][]source}{}{}{}}
\makeatother
\crefname{originalreference}{original reference}{original references}
\Crefname{originalreference}{Original reference}{Original references}
\makeatletter
\newlabel{sites-cohomology-situation-compare}{{0EZ3}{source}{}{sourcepointer-13}{}}
\newlabel{sites-cohomology-situation-compare@cref}{{[originalreference][0][]0EZ3}{[1][0][]source}{}{}{}}
\makeatother
\crefname{originallemma}{original lemma}{original lemmas}
\Crefname{originallemma}{Original lemma}{Original lemmas}
\makeatletter
\newlabel{topologies-lemma-morphism-big-small-etale}{{021I}{source}{}{sourcepointer-14}{}}
\newlabel{topologies-lemma-morphism-big-small-etale@cref}{{[originallemma][0][]021I}{[1][0][]source}{}{}{}}
\makeatother
\crefname{originallemma}{original lemma}{original lemmas}
\Crefname{originallemma}{Original lemma}{Original lemmas}
\makeatletter
\newlabel{topologies-lemma-surjective-proper-ph}{{0DES}{source}{}{sourcepointer-15}{}}
\newlabel{topologies-lemma-surjective-proper-ph@cref}{{[originallemma][0][]0DES}{[1][0][]source}{}{}{}}
\makeatother
\title{The five comparison properties hold for the ph-topology over the small etale site with proper-morphism coefficients}
\author{The Stacks Project Authors}
\date{stacks-project:etale-cohomology.tex:local-snapshot-20261011; GNU FDL 1.2 or later}
\begin{document}
\maketitle
\noindent\textit{Editorial note.} All delivered mathematical text is the author's own, copied from the pinned Stacks Project chapter file etale-cohomology.tex; no statement or proof step was rewritten, reordered or removed.\par\medskip
\noindent\textit{Editorial note.} The counted proof is the author's complete argument: the proof is written out in full: the weak Serre verification with the exactness argument for the extension condition, and the five comparison properties with the reason and reference for each.\par\medskip
\noindent\textit{Editorial note.} The label lemma-compare-ph-etale names the counted statement, and the proof environment immediately following it in the source is delivered verbatim as its proof.\par\medskip
\noindent\textit{Editorial note.} The counted statement is carried by the author's own lemma environment in the chapter source.\par\medskip
\noindent\textit{Editorial note.} Pre-open gate: python operations/agent-authoring/tools/precard\_gate.py --source etale-cohomology.tex --label lemma-compare-ph-etale --start 22321 --end 22386 --local-id 3549 returned RESULT=OK (L3549 free, occupancy FREE). Candidate taken from lane-s4/halfB-priority.json.\par\medskip
\noindent\textit{Editorial note.} The statement and proof are copied from the pinned chapter file; the cited inputs are kept as external pointers with their canonical Stacks tags.\par\medskip
\noindent\textit{Allowed standard prerequisites (editorial).}\begin{itemize}
\item Cohomological descent for the etale comparison. Same chapter, the lemma stating that K -> Ra\_\{X,*\}a\_X\textasciicircum{}\{-1\}K is an isomorphism for torsion coefficients, giving exactness and fullness of pi\_X\textasciicircum{}\{-1\}. Supplies the weak Serre verification and the identification of an extension with a pullback.
\item The description of the pullback in the ph-setting. Same chapter, the lemma identifying pi\_X\textasciicircum{}\{-1\} of a torsion sheaf with the rule given by the sections of its small pullback, in the ph-topology. Supplies property (item-A-sheaf).
\item Base change, torsion and the ph-covering property of proper morphisms. Morphisms, the permanence of properness under base change; Same chapter, the comparison of the small and big higher direct images and the torsion of direct images; Topologies, the lemma that a surjective proper morphism is a ph-covering; More on Morphisms, the domination of an fppf cover by a finite one. Supplies properties (item-base-change-P), (item-A-and-P) and (item-refine-tau-by-P).
\end{itemize}\medskip
\ifdefined\hypertarget\hypertarget{sourcepointer-1}{}\fi
\noindent\textit{Source reference.} Original lemma 0754: sources/stacks-chow/tags.txt, source lines 9118--9118. Prerequisite pointer invoked at proof line(s) 22344; it names a standard fact of a sibling chapter that the counted argument uses, and every step of the counted proof built on it is retained verbatim in the delivered body.\par\medskip
\ifdefined\hypertarget\hypertarget{sourcepointer-2}{}\fi
\noindent\textit{Source reference.} Original reference 0DDI: sources/stacks-etale-cohomology/pinned-original-source.tex, source lines 21308--21308. Prerequisite pointer invoked at proof line(s) 22346; it names a standard fact of the same chapter that the counted argument uses, and every step of the counted proof built on it is retained verbatim in the delivered body.\par\medskip
\ifdefined\hypertarget\hypertarget{sourcepointer-3}{}\fi
\noindent\textit{Source reference.} Original reference 0DDJ: sources/stacks-etale-cohomology/pinned-original-source.tex, source lines 21323--21323. Prerequisite pointer invoked at proof line(s) 22374; it names a standard fact of the same chapter that the counted argument uses, and every step of the counted proof built on it is retained verbatim in the delivered body.\par\medskip
\ifdefined\hypertarget\hypertarget{sourcepointer-4}{}\fi
\noindent\textit{Source reference.} Original reference 0DDW: sources/stacks-etale-cohomology/pinned-original-source.tex, source lines 22175--22175. Prerequisite pointer invoked at proof line(s) 22370; it names a standard fact of the same chapter that the counted argument uses, and every step of the counted proof built on it is retained verbatim in the delivered body.\par\medskip
\ifdefined\hypertarget\hypertarget{sourcepointer-5}{}\fi
\noindent\textit{Source reference.} Original reference 0DDD: sources/stacks-etale-cohomology/pinned-original-source.tex, source lines 13924--13924. Prerequisite pointer invoked at proof line(s) 22377; it names a standard fact of the same chapter that the counted argument uses, and every step of the counted proof built on it is retained verbatim in the delivered body.\par\medskip
\ifdefined\hypertarget\hypertarget{sourcepointer-6}{}\fi
\noindent\textit{Source reference.} Original lemma 0DET: sources/stacks-chow/tags.txt, source lines 16813--16813. Prerequisite pointer invoked at proof line(s) 22382; it names a standard fact of a sibling chapter that the counted argument uses, and every step of the counted proof built on it is retained verbatim in the delivered body.\par\medskip
\ifdefined\hypertarget\hypertarget{sourcepointer-7}{}\fi
\noindent\textit{Source reference.} Original lemma 01W4: sources/stacks-chow/tags.txt, source lines 2469--2469. Prerequisite pointer invoked at proof line(s) 22361; it names a standard fact of a sibling chapter that the counted argument uses, and every step of the counted proof built on it is retained verbatim in the delivered body.\par\medskip
\ifdefined\hypertarget\hypertarget{sourcepointer-8}{}\fi
\noindent\textit{Source reference.} Original reference 0EZ7: sources/stacks-chow/tags.txt, source lines 18717--18717. Prerequisite pointer invoked at proof line(s) 22335; it names a standard fact of a sibling chapter that the counted argument uses, and every step of the counted proof built on it is retained verbatim in the delivered body.\par\medskip
\ifdefined\hypertarget\hypertarget{sourcepointer-9}{}\fi
\noindent\textit{Source reference.} Original reference 0EZ6: sources/stacks-chow/tags.txt, source lines 18716--18716. Prerequisite pointer invoked at proof line(s) 22334; it names a standard fact of a sibling chapter that the counted argument uses, and every step of the counted proof built on it is retained verbatim in the delivered body.\par\medskip
\ifdefined\hypertarget\hypertarget{sourcepointer-10}{}\fi
\noindent\textit{Source reference.} Original reference 0EZ4: sources/stacks-chow/tags.txt, source lines 18714--18714. Prerequisite pointer invoked at proof line(s) 22332; it names a standard fact of a sibling chapter that the counted argument uses, and every step of the counted proof built on it is retained verbatim in the delivered body.\par\medskip
\ifdefined\hypertarget\hypertarget{sourcepointer-11}{}\fi
\noindent\textit{Source reference.} Original reference 0EZ8: sources/stacks-chow/tags.txt, source lines 18718--18718. Prerequisite pointer invoked at proof line(s) 22336; it names a standard fact of a sibling chapter that the counted argument uses, and every step of the counted proof built on it is retained verbatim in the delivered body.\par\medskip
\ifdefined\hypertarget\hypertarget{sourcepointer-12}{}\fi
\noindent\textit{Source reference.} Original reference 0EZ5: sources/stacks-chow/tags.txt, source lines 18715--18715. Prerequisite pointer invoked at proof line(s) 22333; it names a standard fact of a sibling chapter that the counted argument uses, and every step of the counted proof built on it is retained verbatim in the delivered body.\par\medskip
\ifdefined\hypertarget\hypertarget{sourcepointer-13}{}\fi
\noindent\textit{Source reference.} Original reference 0EZ3: sources/stacks-chow/tags.txt, source lines 18713--18713. Prerequisite pointer invoked at proof line(s) 22338; it names a standard fact of a sibling chapter that the counted argument uses, and every step of the counted proof built on it is retained verbatim in the delivered body.\par\medskip
\ifdefined\hypertarget\hypertarget{sourcepointer-14}{}\fi
\noindent\textit{Source reference.} Original lemma 021I: sources/stacks-chow/tags.txt, source lines 2671--2671. Prerequisite pointer invoked at proof line(s) 22366; it names a standard fact of a sibling chapter that the counted argument uses, and every step of the counted proof built on it is retained verbatim in the delivered body.\par\medskip
\ifdefined\hypertarget\hypertarget{sourcepointer-15}{}\fi
\noindent\textit{Source reference.} Original lemma 0DES: sources/stacks-chow/tags.txt, source lines 16812--16812. Prerequisite pointer invoked at proof line(s) 22385; it names a standard fact of a sibling chapter that the counted argument uses, and every step of the counted proof built on it is retained verbatim in the delivered body.\par\medskip
\begin{lemma}
\label{lemma-compare-ph-etale}
Consider the comparison morphism
$\epsilon : (\Sch/S)_{ph} \to (\Sch/S)_\etale$.
Let $\mathcal{P}$ denote the class of proper morphisms of schemes.
For $X$ in $(\Sch/S)_\etale$ denote
$\mathcal{A}'_X \subset \textit{Ab}((\Sch/X)_\etale)$
the full subcategory consisting of sheaves of the form
$\pi_X^{-1}\mathcal{F}$ where $\mathcal{F}$ is a
torsion abelian sheaf on $X_\etale$
Then Cohomology on Sites, Properties
(\ref{sites-cohomology-item-base-change-P}),
(\ref{sites-cohomology-item-restriction-A}),
(\ref{sites-cohomology-item-A-sheaf}),
(\ref{sites-cohomology-item-A-and-P}), and
(\ref{sites-cohomology-item-refine-tau-by-P})
of Cohomology on Sites, Situation
\ref{sites-cohomology-situation-compare} hold.
\end{lemma}

\begin{proof}
We first show that $\mathcal{A}'_X \subset \textit{Ab}((\Sch/X)_\etale)$
is a weak Serre subcategory by checking conditions (1), (2), (3), and (4)
of Homology, Lemma \ref{homology-lemma-characterize-weak-serre-subcategory}.
Parts (1), (2), (3) are immediate as $\pi_X^{-1}$ is exact and
fully faithful for example by Lemma \ref{lemma-cohomological-descent-etale}. If
$0 \to \pi_X^{-1}\mathcal{F} \to \mathcal{G} \to \pi_X^{-1}\mathcal{F}' \to 0$
is a short exact sequence in $\textit{Ab}((\Sch/X)_\etale)$
then $0 \to \mathcal{F} \to \pi_{X, *}\mathcal{G} \to \mathcal{F}' \to 0$
is exact by Lemma \ref{lemma-cohomological-descent-etale}.
In particular we see that $\pi_{X, *}\mathcal{G}$ is an abelian
torsion sheaf on $X_\etale$.
Hence $\mathcal{G} = \pi_X^{-1}\pi_{X, *}\mathcal{G}$ is in
$\mathcal{A}'_X$ which checks the final condition.

\medskip\noindent
Cohomology on Sites, Property (\ref{sites-cohomology-item-base-change-P}) holds
by the existence of fibre products of schemes
and the fact that the base change of a proper morphism of
schemes is a proper morphism of schemes, see
Morphisms, Lemma \ref{morphisms-lemma-base-change-proper}.

\medskip\noindent
Cohomology on Sites, Property (\ref{sites-cohomology-item-restriction-A})
follows from the commutative diagram (3) in
Topologies, Lemma \ref{topologies-lemma-morphism-big-small-etale}.

\medskip\noindent
Cohomology on Sites, Property (\ref{sites-cohomology-item-A-sheaf}) is
Lemma \ref{lemma-describe-pullback-pi-ph}.

\medskip\noindent
Cohomology on Sites, Property (\ref{sites-cohomology-item-A-and-P}) holds by
Lemma \ref{lemma-compare-higher-direct-image-proper} part (2)
and the fact that $R^if_{small}\mathcal{F}$
is torsion if $\mathcal{F}$ is an abelian torsion
sheaf on $X_\etale$, see Lemma \ref{lemma-torsion-direct-image}.

\medskip\noindent
Cohomology on Sites, Property (\ref{sites-cohomology-item-refine-tau-by-P})
follows from More on Morphisms, Lemma
\ref{more-morphisms-lemma-dominate-fppf-etale-locally}
combined with the fact that a finite morphism is proper
and a surjective proper morphism defines a ph covering, see
Topologies, Lemma \ref{topologies-lemma-surjective-proper-ph}.
\end{proof}

\end{document}
